\chapter{マシンのデバッグ}
\chaptermark{マシンのデバッグ}
\label{sec:debugging}

\section{回路図のないマシンをどうデバッグするか}

回路図なしで動作中の基板を渡されたエンジニアは、四つの手法でデバッグする。\emph{プローブ}を当てて配線に何が流れているかを見る。\emph{テスト信号を入れて}何が変わったかを見る。回路の一部を\emph{切り離して}何が動かなくなったかを見る。そして\emph{制御レジスタを読んで}、基板がどのモードで動いているかを知る。本書全体が脳の回路図を読む試みである。本章では、エンジニアが四つの手法のうちどれをすでに使えるのか、そして前章までの内容から何が期待できるのかを見る。

今日このようなデバッグの最も目立つ例は、ブレイン・コンピュータ・インターフェースである。ニューロンのスパイクを読み取って外部の機械への指令に変える装置、あるいは逆に外から脳へ信号を入れる装置である。本章の大部分はこれを扱い、とくにNeuralink社の取り組みを取り上げる。

\section{プローブ：電極には何が聞こえるか}

脳のプローブは、組織に刺し込んだ細い電極である。電極には、半径およそ100マイクロメートル以内にあるニューロン、ふつうは1個から数個のスパイクが聞こえる。電極上のスパイクは、数十から数百マイクロボルト、幅約1ミリ秒のパルスであり、その波形から隣り合うニューロンを区別できる。最もよく聞こえるのは大きな\nt{GLUT}ニューロンである。皮質ではこれが多数派であり（表~\ref{tab:types}）、電流も強い。

すぐに主要な困難が見えてくる。今日の優れたプローブは数百から数千の電極をもつが、脳のニューロンは$8.6\cdot10^{10}$個ある。盗聴しているのはマシンのおよそ1億分の1にすぎない。これほどわずかな標本から、なぜ何かがわかるのか。答えは第\ref{sec:code}章にある。隣り合うニューロンのノイズは相関しており、群での平均化には限界がある（命題~\ref{prop:pooling}）。100個のニューロンは1000個とほぼ同じだけの情報を運ぶ。さらに、運動皮質が解く課題は低次元である。腕の関節は多くなく（第\ref{sec:muscles}章）、腕を制御する数百個のニューロンの活動は、わずか10--20個の共通変数の空間に収まる~\cite{Gallego2017}。100ニューロンのプローブは、これらの変数のほぼすべてを聞き取れる。もし脳が各ニューロンに独立したメッセージを格納していたら、このような盗聴は役に立たなかっただろう。

\section{データストリーム：プローブ上での圧縮}

プローブが出すデータは膨大である。スパイクの波形を見るには、各電極を毎秒約$2$万回、約10ビットの精度でデジタル化する必要がある。1000個の電極では
\begin{empheq}[box=\keybox]{equation}
\begin{aligned}
R_{\text{raw}} \;&=\; N\,f_s\,b \;\approx\; 10^3\cdot2\cdot10^4\cdot10 \;\approx\; 2\cdot10^8\ \text{bit/s},\\
R_{\text{spk}} \;&\approx\; N\,r\,\log_2\frac{e}{r\,\Delta t} \;\approx\; 8\cdot10^4\ \text{bit/s}.
\end{aligned}
\label{eq:probedata}
\end{empheq}
となる。第2の見積もりは、$r=10$スパイク/s、精度$\Delta t=1\ms$でのスパイク列の容量~\eqref{eq:spikeentropy}である。スパイク列に実際に含まれるものすべてが、2500分の1の量に収まる。埋め込み型装置にとってこれは決定的な違いである。頭の中で許される電力では、組織を加熱せずに生のデータを皮膚越しに無線で送ることはできない。したがって埋め込み型装置は、電極のそばのチップ上でスパイクを検出し、イベント、すなわちチャネル番号とスパイクの時刻だけを外に送る必要がある。Neuralinkのシステムの最初の報告では、まだそこまで到達していなかった。チップは信号を増幅してデジタル化し、生のデータはすべてケーブルで外へ送られ、スパイクは外部のプログラムが検出していた。チップ上の圧縮と無線接続は計画として挙げられていた~\cite{Musk2019}。

これはおなじみの手法である。眼は長いケーブルに送る前に信号を100分の1に圧縮する（第\ref{sec:sensors}章）。イベントカメラはフレームではなく変化を送る。配線に収めるために、プローブは脳自身と同じことをせざるをえない。スパイクで、しかも新しいことだけを伝えるのである。

\section{Neuralinkは何をしているか}

Neuralinkの装置は、これらの制約に合わせて組み立てたプローブである~\cite{Musk2019}。硬い針の配列の代わりに、髪の毛より細い柔軟なスレッドを多数使い、各スレッドに沿って電極を並べる。最初の報告では$96$本のスレッドに最大$3072$個の電極、1本あたり$32$個であった。同社によれば、ヒトに埋め込んだ装置は$64$本のスレッドで約1000チャネルである。柔軟なスレッドは組織の損傷が少なく、頭蓋内で脳が常にわずかに動くことにもよく耐える。スレッドはロボットが血管を避けながら挿入する。前述のとおり、最初の報告では生のデータ~\eqref{eq:probedata}はケーブルで外へ送られていた。同社によれば、ヒト用の装置は別の構成である。筐体は完全に皮膚の下にあり、スパイクは内部で検出され、データは無線で送られ、電力はワイヤレスで供給される。

最初の装置の目的は読み取りである。スレッドを腕の運動を計画する皮質領域に置き、デコーダがスパイクを画面上のカーソルの動きに変える。腕が麻痺した人が、動きを想像してコンピュータを操作する。発想自体は新しくない。100電極の硬い配列を使うインターフェースは、以前からカーソル操作~\cite{Hochberg2006}、手書きの動きを想像して文字を書くこと~\cite{Willett2021}、さらには話そうとする試みを毎分約60語の速さで画面上のテキストに変換すること~\cite{Willett2023}を可能にしてきた。Neuralinkの新しさはエンジニアリングにある。チャネル数、ワイヤレス化、密閉筐体、ロボットによる挿入、すなわち研究室の外で何年も装着できるプローブを作る試みである。我々の知る限り、ヒトでの試験結果は主に同社自身が公表しており、査読付き論文ではない。同社は特に、挿入後に一部のスレッドがずれ、動作するチャネルの数が減ったと報告している。デバッグではよくある不具合である。基板に対して動くプローブは、もう別の配線を聞いている。

同社の第二の方向は書き込みである。視力を失った人の視覚皮質に信号を入れる視覚用の装置である。これについては後の書き込みの節で述べる。

\section{読み取り：受け手としてのデコーダ}

インターフェースのデコーダは、命題~\ref{prop:reader}の意味での受け手である。メッセージのどの部分を読むかは自分で決める。実際には、ほとんどすべてのインターフェースが\emph{群の発火率}を読む。各チャネルのスパイクを数十ミリ秒の窓で数え、意図した動きが得られるように選んだ重みで足し合わせる。これは第\ref{sec:rate}節の群のレート符号であり、偶然選ばれたのではない。窓あたり数スパイクでは1個のニューロンの発火率は決まらないが、100個の和ならもう使える。

どれだけの情報を取り出せるか。「想像上の手」による書字は毎分約90文字であった~\cite{Willett2021}。つまり毎秒1.5文字であり、1文字5ビットとしても毎秒8ビット未満である。Neuralinkによれば、カーソル操作は毎秒約10ビットを与える。一方、上限~\eqref{eq:spikeentropy}は\emph{1個}のニューロンで毎秒数十ビット、100個なら数千ビットである。インターフェースは、盗聴したニューロンが運べるはずの情報のごく一部しか取り出していない。ボトルネックは配線ではない。群が運ぶ独立変数の数（命題~\ref{prop:pooling}）と、第\ref{sec:code}章で見たように完全には解読されていない符号をデコーダがどれだけ理解しているかである。

第\ref{sec:motorlearning}章でおなじみの第二の特徴もある。デコーダと脳は互いから学ぶ。脳はカーソルがどこへ動いたかを見て誤差を受け取り、指令を調整する。これは誤差配線による運動学習と同じである。一方エンジニアは、スレッドの下のニューロンが変わり、ネットワークの結合が学習し直されるので、ときどきデコーダの重みを選び直す。こうして互いに合わせ合う二人の学習者ができる。この組が暴走しないようにするには、学習速度を分けておくとよい。一方は速く、他方は遅く学ぶようにする。

\begin{keyblock}
\begin{proposition}[1000ニューロンのプローブはなぜ働くか]
\label{prop:probe}
ニューロンのごく一部を盗聴するだけのインターフェースが動きを制御できるのは、群の活動が低次元で、そのノイズが相関しているからである。数百個のニューロンが共通変数のほぼすべてを運ぶ。同じ理由で、インターフェースが取り出すのは毎秒数ビットにすぎない。これはスパイク列の容量ではなく、課題の独立変数の数で決まる。インターフェースの性能は測定済みである。符号が理解されたときデコーダがどれだけ良くなるかは未解決の問題である。
\end{proposition}
\end{keyblock}

\section{書き込み：読み取りより難しい}

マシンに信号を入れるのは、読むより難しい。電流を流す電極は1個のニューロンではなく、周囲のすべてのニューロンと配線を、種類も符号も区別せずに興奮させる。\emph{どの}ニューロンが\emph{いつ}発火したかが重要な符号（第\ref{sec:code}章）を、これで書き込むことはできない。大まかに\emph{どこで}、\emph{どれだけ強く}を指定できるだけである。

視覚にはこれで部分的に足りる。視覚皮質の電極に電流を流すと、人は視野の特定の位置に光る点を見る。最大16個の点を同時に点灯させると、視力を失った人がいくつかの文字を認識し、物体の境界を区別できた~\cite{Fernandez2021}。これは場所による符号であり、書き込むことができる。しかし第\ref{sec:sensors}章を思い出してほしい。眼が脳に送るのはピクセルではなく、圧縮された記述である。エッジ、変化、明るくなることと暗くなることが、別々の配線で送られる。皮質に「光の点」を書き込むのは、入力にまったく別の符号を待っているマシンにピクセルを書き込むことである。だから人工視覚は、電極の数から期待されるより今のところ粗い。電極の要らないテスト信号もある。錯視は眼を通して普通でない入力を与え、マシンを通常のモードから外す。そして誤りの一つ一つが、それぞれの機構を指し示す（第\ref{sec:illusions}節）。

\section{プローブに見えないもの}

電極が聞くのはアドレスバス、つまり\nt{GLUT}ニューロンのスパイクと、それより聞き取りにくい小さな\nt{GABA}ニューロンのスパイクである。しかし第\ref{sec:serocomputer}--\ref{sec:acet}章で見たように、マシン全体の動作モードはブロードキャストレジスタ、すなわち\nt{SERO}、\nt{DOPA}、\nt{NORA}、\nt{ACET}の濃度が決める。電極にはこれが聞こえない。発信源のニューロンはごくわずかであり、その信号はプローブのそばのスパイクではなく、体積中の濃度だからである。データバスは見えても制御レジスタが見えないデバッガは、いつまでも驚き続けるだろう。同じ入力が異なる応答を生む。どこかで$g$、$a$、$\tau_\gamma$が変わったからである。濃度は組織内の化学センサで直接測れるが~\cite{Bunin1998}、インターフェースではそうしたセンサはまだ使われていない。マシンの第二の遅い出力、すなわち血中のホルモン（第\ref{sec:chemoutput}章）も、プローブの視野のまったく外にある。これは電極ではなく採血で測る。

プローブには重みも見えない。そして記憶と学習したことのほぼすべては、まさに重みに格納されている。重みは応答の変わり方から推測するしかない。また、プローブには人が感じるような痛みも見えない。第\ref{sec:pain}章で見たように、警報信号の強さはマシン自身、すなわち脊髄のゲートと上位の調節器が決めるので、センサからどれほど痛いかは読み取れない。

\section{まとめ：デバッグと未解決の問題}

\begin{table}[t]
\centering
\footnotesize
\setlength{\tabcolsep}{4pt}
\renewcommand{\arraystretch}{1.15}
\begin{tabular}{>{\raggedright}p{2.2cm}>{\raggedright}p{3.6cm}>{\raggedright}p{3.4cm}>{\raggedright\arraybackslash}p{4.0cm}}
\toprule
デバッグ手法 & インターフェースがすること & 本書による説明 & わかっていること \\
\midrule
プローブ & 数百--数千個のニューロンを聞く & 低次元性と相関したノイズ（第\ref{sec:code}章） & 測定済み \\
プローブ上での圧縮 & 生の信号の代わりにイベントを送る & スパイク列の容量~\eqref{eq:spikeentropy} & 工学的に解決済み \\
読み取り & 群の発火率$\to$動き、文字、発話 & デコーダは受け手、群のレート符号 & 動作する。毎秒数ビット \\
共同学習 & 脳とデコーダが互いに合わせ合う & 誤差配線（第\ref{sec:motorlearning}章） & 観察されている。調整の規則は研究課題 \\
書き込み & 電流が視野に点を灯す & 場所の符号は書き込めるが、時間の符号は書けない（第\ref{sec:sensors}章） & 点は測定済み。実用的な視覚はまだ \\
レジスタの読み取り & ほとんど行われていない & ブロードキャストチャネル（第\ref{sec:serocomputer}--\ref{sec:acet}章） & 化学センサは実験段階 \\
\bottomrule
\end{tabular}
\caption{マシンのデバッグ：ブレイン・コンピュータ・インターフェースにできることと、本書の各章がそれについて語ること。}
\label{tab:debugging}
\end{table}

表~\ref{tab:debugging}が本章をまとめる。ブレイン・コンピュータ・インターフェースは、すでにアドレスバスにプローブを当て、毎秒数ビットを読み取れる。それがそもそも働くのは、第\ref{sec:code}章の低次元性と相関したノイズのおかげである。マシンへの書き込みは粗くしかできない。マシンの入力符号は圧縮され、個々の配線にアドレス付けされているからである。そして制御レジスタと重み、すなわちマシンを学習可能で調整可能にしているものは、プローブにはまだほとんど見えない。

第\ref{sec:synthesis}章の未解決の問題は、どれもデバッガへの課題でもある。どの符号を読むべきかは、プローブがより高密度になり、デコーダが発火率だけでなくスパイクのタイミングを使えるようになったときに決着する。多層ネットワークがどう学習するかは、複数の層に同時にプローブを置き、学習中に応答がどう変わるかを見れば確かめられる。ブロードキャストチャネルが何を伝えているかは、電気的なプローブに化学的なプローブが加わったときに見えてくる。本書は、マシンの回路図をその振る舞いと、どのエンジニアも知っている法則から読み取る試みである。いま作られているデバッガが、この回路図をプローブで確かめることを可能にするだろう。

\section{問題}

\begin{exercise}[\lvlA]
\label{ex:17:1}
\emph{無線チャネルの予算。}皮膚越しの伝送は1ビットあたり$1$~nJかかり、送信機に使える電力は$10$~mW以下である。$N=1000$電極の生データ~\eqref{eq:probedata}とスパイク列を送るのに必要な電力はそれぞれいくらか。生データなら1ビットあたり何ジュールに抑える必要があるか。
\end{exercise}

\begin{exercise}[\lvlA]
\label{ex:17:2}
\emph{100ニューロンに何ビットあるか。}デコーダは$N$個のニューロンのスパイクを$50\ms$の窓で足し合わせる。$r=20$~スパイク/s、ノイズの対相関$c=0.1$、信号は群の発火率を（二乗平均で）$30\%$変える。$N=10$、$100$、$1000$、$N\to\infty$での窓あたりのレート$\tfrac12\log_2(1+\text{SNR})$を見積もれ。$c=0$、$N=1000$ではどうか。
\end{exercise}

\begin{exercise}[\lvlB]
\label{ex:17:3}
\emph{キーボード型インターフェースの速度。}インターフェースは毎秒1.5回、$N=32$文字から1文字を選ぶ。意図した文字を選ぶ確率は$P$、誤りは等確率とする。$P=0.99$、$0.94$、$0.8$で毎秒何ビット伝えるか。$P=0.94$で$10$~bit/sを得るには毎秒何文字必要か。
\end{exercise}

\begin{exercise}[\lvlB]
\label{ex:17:4}
\emph{シミュレーション：群のデコーダ。}$N$個のニューロンで$\lambda_i=[b+m\,\mathbf v\cdot\mathbf p_i]_+$、$b=20$、$m=15$~スパイク/s、窓$50\ms$、$\mathbf v$の各成分の標準偏差$0.5$とする。全ニューロンが共通ノイズ入りの$\mathbf v+\boldsymbol\xi$、$\sigma_\xi=0.2$を見る。群ベクトルによる不偏デコーダをシミュレートせよ。二つのノイズが等しくなる$N$はいくつか。$N\to\infty$で窓あたり成分あたり何ビット得られるか。
\end{exercise}

\begin{exercise}[\lvlB]
\label{ex:17:5}
\emph{二人の学習者。}脳は指令$m=b\,v^*$を、デコーダは$v=d\,m$を出す。$\langle v^{*2}\rangle=1$。両者は試行ごとに$\tfrac12(v-v^*)^2$について学習率$\eta$で勾配を1ステップ下る。$b=1.2$、$d=1$から始め、$\eta=0.8$、$1.1$、$1.5$、$2$でシミュレートせよ。それぞれ単独なら$\eta<2$で安定である。組が安定なのは$\eta$がいくつのときか。
\end{exercise}

\begin{exercise}[\lvlB]
\label{ex:17:6}
\emph{設計：プローブのパイプライン。}$1024$チャネル、毎秒$2$万サンプル、ニューロンは$10$~スパイク/s、無線は$1$~nJ/bitとする。サンプルの白色ノイズに対してどの閾値なら、偽スパイクがチャネルあたり$0.1$~Hz以下になるか。$20\ms$ごとのスパイク数を$2$ビットで送るとき、電力はいくらで、$10$ビットの生サンプルの何分の1か。
\end{exercise}

\begin{exercise}[\lvlC]
\label{ex:17:7}
\emph{インターフェースの速度限界。}帯域$W=5$~Hzの$D=10$変数について、$\text{SNR}\approx0.9$（信号に似たノイズ）と$90$（$1000$ニューロンの独立ノイズ）の場合に$R\le DW\log_2(1+\text{SNR})$を見積もれ。測定値は約$10$~bit/sである。このギャップの二つの説明、信号に似たノイズと符号の無知とを区別する実験は何か。
\end{exercise}
